\documentclass[11pt]{article}
\usepackage[margin=1in]{geometry}
\usepackage[T1]{fontenc}
\usepackage[utf8]{inputenc}
\usepackage{lmodern}
\usepackage{hyperref}
\usepackage{xcolor}
\usepackage{fancyvrb}
\usepackage{listings}
\usepackage{CJKutf8}
\lstset{basicstyle=\ttfamily\small,columns=fullflexible,breaklines=true,upquote=true,frame=single,framerule=0.2pt,tabsize=2}

\title{\texttt{\detokenize{model/egmn.py}} (Q\&A Documentation)}
\author{}
\date{}
\begin{document}
\begin{CJK*}{UTF8}{gbsn}
\maketitle

\paragraph{Q: What does this file implement at a high level?}
\textbf{A:} This file implements EGMN: a mixture-distribution watch-time predictor. It encodes sparse/sequence/continuous features, produces mixture parameters (mixture logits, exponential rate, Gaussian means/variances), and provides both a likelihood-based loss and a scalar prediction via the mixture mean.

\paragraph{Q: Why is a mixture distribution needed for watch-time?}
\textbf{A:} Watch-time has a heavy spike near 0 (quick-skip) and a long tail. A single regression head struggles to represent both. The model uses an Exponential component for the spike and multiple (truncated) Gaussians to represent the tail/modes.

\paragraph{Q: What are the model outputs and how are they used?}
\textbf{A:} Forward returns (pi, lambda, mu, sigma). Training uses them in the negative log-likelihood and auxiliary terms; inference uses them to compute a scalar expectation (mixture mean) as the final predicted watch-time.

\paragraph{Q: What does this file export (public API)?}
\textbf{A:} The public API is the set of classes/functions other modules import (sometimes re-exported by package initializers). The key top-level definitions detected here are:

\begin{Verbatim}[fontsize=\small]
class EGMN  # methods: __init__, build, init, forward, loss, get_quantile, predict
\end{Verbatim}

\paragraph{Q: What are the main dependencies (imports) and why do they matter?}
\textbf{A:} Imports show whether this file is model code (torch), data code (pandas), or evaluation/analysis (sklearn). They also determine runtime requirements.

\vspace{1mm}
\VerbatimInput[firstline=1,lastline=50,fontsize=\small]{/workspace/model/egmn.py}


\paragraph{Q: What is \texttt{\detokenize{EGMN}} and what responsibility does it have?}
\textbf{A:} This class is a core unit in the pipeline. Its responsibility is inferred from its name and how run scripts call it.

\paragraph{Q: What are the important methods on \texttt{\detokenize{EGMN}}?}
\textbf{A:} In this repo, methods like forward/loss/predict define computation and training targets. The methods found are:

\begin{Verbatim}[fontsize=\small]
__init__
build
init
forward
loss
get_quantile
predict
\end{Verbatim}

\paragraph{Q: What does the implementation look like for \texttt{\detokenize{EGMN}}?}
\textbf{A:} Excerpt from the file:

\vspace{1mm}
\VerbatimInput[firstline=9,lastline=89,fontsize=\scriptsize]{/workspace/model/egmn.py}


\paragraph{Q: Why is \texttt{\detokenize{EGMN}} needed by training scripts?}
\textbf{A:} Training scripts assume this class can be instantiated from the dataset description and can map a feature dictionary to outputs required by that script's loss. This separation makes it easy to swap models while reusing the same dataloader and metrics.

\paragraph{Q: Example: end-to-end data flow involving this file}
\textbf{A:} Core pipeline shape (schematic):

\begin{Verbatim}[fontsize=\small]
for features, label in dataloaders["train"]:
    out = model(features)
    loss = ...
    loss.backward(); optimizer.step()
\end{Verbatim}

\paragraph{Q: Full source code of the Python file}
\textbf{A:} The complete file is included verbatim for reference.

\VerbatimInput[fontsize=\scriptsize]{/workspace/model/egmn.py}

\end{CJK*}
\end{document}